\documentclass[conference]{IEEEtran}
\usepackage{amsmath,amssymb}
\usepackage{booktabs}
\usepackage{url}
\usepackage[hidelinks]{hyperref}

\begin{document}

\title{Six Ways a Network Simulation Study Produces a Wrong Number \\
\large A self-audit of one LoRa/802.11 co-design, with the guards that caught each}

\author{\IEEEauthorblockN{Mohamed A. Farouk}
\IEEEauthorblockA{AUTHBC Project --- methodological companion (auto-generated from frozen results)}}

\maketitle

\begin{abstract}
Twenty years ago Kurkowski \emph{et al.} found that none of 84 MobiHoc simulation papers referenced
public code, and follow-up work reports that fewer than 15\% of MANET simulation studies are
reproducible. Those surveys tell us the literature has a credibility problem; they do not tell an
author what, mechanically, goes wrong. We report the other half: a self-audit of a single
LoRa/802.11 authentication co-design in which \textbf{ten quantitative results moved or were
retracted}, and a taxonomy of the \textbf{six distinct mechanisms} responsible. Only one is fixed
by running more seeds. Two were protected by a passing test suite. Three were internal
contradictions between a paper and its own artifacts that every gate in the project reported green.
For each class we give the concrete instance, the number it moved, and the mechanical check that now
prevents it --- a bootstrap over threshold crossings, an RNG-stream isolation assertion, a test that
parses the paper's \LaTeX{} tables and compares them to the CSVs, and a test that compares the
documentation's own counts to reality. We argue that \emph{seed discipline is necessary and badly
insufficient}, and that the useful unit of reproducibility work is not the artifact but the
\emph{invariant a machine can check}.
\end{abstract}

\section{Introduction}
Simulation credibility in networking has been studied as a \emph{reporting} problem. Kurkowski,
Camp and Colagrosso surveyed MobiHoc 2000--2004 and found that experiments were rarely repeatable,
unbiased, rigorous and statistically sound at once~\cite{kurkowski2005incredibles}; a decade later
the VANET community reached similar conclusions~\cite{ccr2018vanet}. The prescription that followed
is familiar: publish code, report seeds, give confidence intervals.

We took that prescription seriously, and it was not enough. Over a five-week audit of one
co-design study --- an authenticated UAV telemetry ledger evaluated with NS-3 and two Raspberry~Pi
radios --- ten quantitative results moved or were withdrawn. Every one had passed the project's own
gates. \textbf{Seed discipline caught four of them. It was structurally incapable of catching the
other six.} Four later review passes over the write-up moved roughly thirty further numbers, and
the dominant mechanism in those was not any of the five: it was corrections that had been made
correctly at the source and never propagated (C6).

This paper is the taxonomy. It is deliberately a self-audit: every instance is our own, with the
before and after numbers, because a defect list assembled from other people's papers cannot say what
the wrong number \emph{was}.

\section{The five classes}

\subsection{C1 --- A small-sample mean compared against a threshold}
The known one. A capacity figure is defined as the largest $N$ whose mean delivery still meets
$V \geq 0.95$; with three seeds the mean lands on the wrong side of the threshold.
\textbf{Instance:} $N_{\max}$ fell from 5 to 3 at 30 seeds; a delay crossing moved
$U{=}2.797 \to 2.435$; a validation band endpoint moved $-1.44\% \to -0.51\%$.
\textbf{Guard:} 30 seeds by default, and drivers that emit min/max/$\sigma$.
\textbf{This is the only class more seeds fix.}

\subsection{C2 --- An unverified constant on the measurement path}
\textbf{Instance:} a two-node broadcast experiment returned $97.45\%$ delivery with
$\sigma = 0.196$\,pp over eight windows --- tight, repeatable, and meaningless. 802.11 sends
broadcast at the lowest basic rate, which on 2.4\,GHz is 802.11b's 1\,Mb/s; a 1400\,B frame then
occupies 11.2\,ms, so the nominal ``100\,fps operating point'' was $118\%$ of capacity. The
measurement was saturation, not channel loss.
\textbf{Why seeds cannot help:} eight windows agreed to 0.2\,pp \emph{on the wrong quantity}.
Repeatability is not validity.
\textbf{Guard:} state the assumption the expectation depends on, and put a discriminator in the same
run --- an offered-load sweep beside the point measurement located the ceiling at 501\,fps and
exposed it immediately.

\subsection{C3 --- A threshold applied to a mean rather than to a distribution}
\textbf{Instance:} at the certified $N{=}3$, \emph{nine of thirty} runs failed the very criterion
being certified. The mean cleared $0.95$; the distribution straddled it. Reported as a point,
$N_{\max}{=}3$ has a 95\% bootstrap CI of $[2,3]$; read per realisation --- at least 95\% of runs
must individually meet $V$ --- it is \textbf{1}.
\textbf{Why it is not C1:} the sample was already 30. The defect is that a reliability target was
evaluated as a throughput average.
\textbf{Guard:} a nonparametric bootstrap over the crossing itself, since a crossing is discrete and
does not inherit a mean's interval.

\subsection{C4 --- A configuration change that perturbs the random realisation}
\textbf{Instance:} installing a mobility model appeared to cost 5 percentage points of delivery.
NS-3 assigns RNG stream indices in object-creation order, and mobility models are constructed before
senders, so the comparison arms differed in bookkeeping rather than physics. The tell was that
5\,m/s and 20\,m/s gave \emph{identical} delivery --- a positional effect would scale with
displacement.
\textbf{Guard:} pin the senders' streams by node id, and assert an invariant the physics forces:
\texttt{sent} counts transmissions, so no propagation or gateway option may change it. Applied
across the study's other comparison axes, that assertion found them clean.

\subsection{C5 --- A claim wider than the experiment supports}
\textbf{Instance:} a two-node hardware measurement with one transmitter validates airtime and link
loss. It cannot validate a contention model, which describes $N$ competing stations --- and an
earlier draft implied it did. Separately, a figure was quoted from the correct paper but the wrong
figure, inverting a comparison.
\textbf{Guard:} the discipline of writing what a result cannot show beside what it can, and
\textbf{quoting the figure, not the paper}.

\subsection{C6 --- A correction that does not propagate}
The class we did not anticipate, and the one with the most instances. C1--C5 describe how a number
is \emph{produced} wrongly. C6 is different: the number is corrected \emph{once}, at its source,
and the correction fails to reach every place the value already appears. The source of truth and the
artifacts agree with each other; only the derived copies disagree, and they disagree silently.

\textbf{Instances}, all from four further review passes over a single paper:
\begin{itemize}\itemsep2pt
\item a frame header measured at 44\,B replaced an assumed 40\,B --- and \textbf{five constants
\emph{derived} from it kept the old value}. Searching for ``40'' finds none of them, because none of
them prints the header: they print $(M{-}H_f{-}g_a)/(b{+}1)$, $M/(M{-}H_f{-}g_a)$, a batch ceiling.
\item a finding was \textbf{retracted} in the audit register and remained printed in a table, one
hundred lines below a sentence asserting its opposite in bold;
\item six three-seed artifacts were \textbf{deleted} as superseded, and a table still built on one
of them matched the deleted file to three decimals;
\item a threshold crossing corrected in a table caption was left uncorrected on a figure's axis;
\item a simulator upgrade documented in the paper's own limitations was left stale in a plot title;
\item a completed phase left a figure footing ``pending''.
\end{itemize}

\textbf{Why seeds, artifacts and code tests are all blind to it.} There is no wrong computation to
find. Every artifact is correct and mutually consistent; the defect lives in the copies --- prose,
captions, axis labels, plot titles --- which no gate reads. Two of the six had \emph{already been
found and fixed elsewhere in the same repository}, which is what makes the class distinctive: the
project had the right answer written down and shipped the wrong one anyway.

\textbf{Guard.} Re-derive every consumer when an input changes, and treat a retraction as a
\texttt{grep} obligation rather than a register entry. Mechanically: parse the prose for the
constants it quotes and recompute them from the measured inputs. Note the asymmetry --- a
half-correct table is \emph{harder} to catch than a wholly wrong one, because the columns that
agree with their source make the table read as verified.

\section{What the gates missed, and why}
Three of the ten corrections were \emph{internal contradictions} in which a paper disagreed with its
own artifacts while every test passed:

\begin{itemize}
\item a capacity table kept superseded values ($233/116$) after the prose was corrected to
$213/100$;
\item a sentence stating a ratio ``is stable under either threshold'' sat in the same paper as a
sentence warning that exactly that rounding is misleading ($3.22\times$ vs $2.42\times$);
\item the abstract filed the study's most durable result under ``applications of established
results'' while the conclusion called it ``the more durable contribution''.
\end{itemize}

None was a code defect, so no code test could see it. \textbf{The unit of reproducibility work that
would have caught them is not an artifact but a machine-checkable invariant across the boundary
between prose and data.} We now parse the paper's \LaTeX{} tables and compare every cell to the
generating CSV, and parse the project documentation for asserted counts and compare them to reality.
The second test found four stale claims on its first run.

Two defects were \emph{protected by passing tests}. A unit test asserted
\texttt{channel\_utilisation(1) == 0.0} with the comment ``a lone sender never contends'' --- true of
contention, false of utilisation, and it had encoded the defect as intended behaviour. A green suite
means the code does what the tests say, not that the tests say the right thing.

\section{Pre-registration is cheap and it works}
The study's success criterion was committed two days before the result it judges, and both commits
are public, so the ordering is checkable by a third party rather than asserted. A second
pre-registration claim, on a later sub-study, had to be \textbf{withdrawn}: it cited an expectations
file that had never been committed. We did not reconstruct it. Writing that file after the outcome
was known would have manufactured evidence that is unfalsifiable from outside.

When a follow-up survey was later planned, its protocol --- hypothesis, inclusion criteria,
falsification threshold, keyword set, and the rule that unreadable sources are excluded from the
denominator rather than scored as negatives --- was committed in a data-free commit before the
corpus existed. \textbf{It then cost us the claim, which is the strongest evidence we have that it
worked.}

The hypothesis was that ns-3 LoRa simulation studies do not report the replication information a
reader needs. The threshold for abandoning it --- $25\%$ of the corpus reporting --- was fixed in
advance. As retrieval improved the estimate walked steadily toward that line: $0\%$ at $n{=}4$,
then $14.3\%$, $15.0\%$, and $21.7\%$ at $n{=}23$, with a $95\%$ interval of $[7.5, 43.7]\%$.
\textbf{The interval contains the threshold.} The test could no longer distinguish its hypothesis
from its own falsification, so the claim was withdrawn from the paper and the corpus kept in the
repository as a null result.

Two details are worth stating because both were temptations. The point estimate, $21.7\%$, is on
the favourable side of $25\%$ and could have been quoted alone. And the subset of the corpus not
shaped by what a script could download --- the properly published papers --- reads $28.6\%$,
\emph{above} the threshold, a split we had flagged as a confound two rounds earlier. It is not
significant ($p = 0.61$) and is not reported as a finding, because a non-significant subgroup
discovered after the fact is the forking-paths error, not a result. \textbf{Pre-registration is
cheap; what makes it work is honouring it when the number lands on the wrong side.} Within an hour it forced two corrections to our own claims: a phenomenon we had
called unobserved turned out to have prior art, and a tally of ``nine of nine'' fell to four under
the criteria we had just written down. \textbf{Both would otherwise have reached a submission.}

\section{Conclusion}
The credibility literature tells authors to publish code and report seeds. Both are necessary. In
one audited study, seeds addressed four defects out of ten; the rest came from an unverified
constant, a criterion applied to the wrong statistic, a comparison confounded by RNG bookkeeping, an
overreaching claim, and three contradictions between a paper and its own data. Four further passes
over the write-up then found a sixth and larger class, in which the study had \emph{already made}
the correction and shipped the uncorrected copy regardless.

The transferable lesson is not a longer checklist. It is that \emph{every claim worth making should
have an invariant a machine can check}, including the claims a paper makes about itself --- and that
the most valuable checks are the ones that can return an answer you did not want. Ours did, ten
times.

\begin{thebibliography}{9}
\bibitem{kurkowski2005incredibles} S.~Kurkowski, T.~Camp, and M.~Colagrosso, ``MANET simulation
studies: the incredibles,'' \emph{ACM SIGMOBILE Mobile Computing and Communications Review},
vol.~9, no.~4, pp.~50--61, 2005.
\bibitem{ccr2018vanet} ``VANETs' research over the past decade: overview, credibility, and trends,''
\emph{ACM SIGCOMM Computer Communication Review}, 2018.
\end{thebibliography}

\end{document}
